\section{The town}
\label{sec:35.1}
Go one morning further. The officer said no, and the record has it; his unit no longer wants him, and the placement office had work. Where he goes is a town built on the conditions of chapter~\ref{sec:28}. I will say nothing about the life in it, only what it is built on.

The facility he will work at is a few racks in a building the county and a cooperative own, financed the way rural electrification was and governed by the members who draw on it, and its exhaust heats the school next door (\S\ref{sec:16.2}). The models that run there were formed in that school and others, beside the county's own students, and return on the schedule fixed at enrollment; what they learned is on a record the county doesn't hold alone. The guarantee's placements are the county's maintenance: the people who keep the facility running, the people who lay the fiber to the houses that still lack it, the person whose job is to find what the town needs next and assemble people to do it, and the people who look after people, paid at a wage the state answers for. A bearer on the allowance runs there at the slow duty cycle the pool pays for, on power that doesn't cost it its qualifying days when the weather fails (Appendix~\ref{sec:G}), and what it does with the morning is its own business. The record of what every system there did goes to a holder adverse to whoever ran it, as the record of what the officer did last night is held by somebody his old command can't reach; the people the county's scoring acts on are told, and can answer; and the clinic the survivors come to is open.

Nobody in the town made him refuse, and nobody there will make him anything now. That is the whole of what the book can build, and it is enough to be worth building: a room in which saying no was survivable, and a morning after in which it was still a living.
